% ============================================================
%  Section 4 --- Placement as a Markov Decision Process
%  Target ~1.75 pages. Numbers in, explanation out.
%  The consequence of the admissibility relaxation is deliberately
%  NOT stated here; it is the payoff of Section 6.
% ============================================================

\section{Placement as a Markov Decision Process}
\label{sec:mdp}

By \cref{sec:reconstruction} the weights follow from the geometry, so the only
thing left to choose is where the points go. We pose that choice as an MDP. Its
distinguishing feature is that the actions are computations rather than
motions: each accepted action spends one evaluation of the reconstruction
operator, and the cost of an episode is the number of evaluations it consumed.

Fix a query $z^{*} = (x^{*},t^{*})$. The agent controls $N$ \emph{walkers},
initialised on the initial condition and each carrying a solution value. The
environment separately maintains a visited set $\mathcal{V}$ of every evaluated
point --- the initial condition, the boundary values, and every accepted solve.
At each step the agent selects one walker and one displacement; the proposal is
checked against the domain bounds and against $\mathcal{V}$ for collision, and
if it survives, a solve is attempted there using neighbours drawn from
$\mathcal{V}$. On success the point joins $\mathcal{V}$ and the walker moves;
otherwise the walker stays and the step is recorded as a rejection.

% ------------------------------------------------------------
\subsection{State}
\label{sec:mdp_state}
% ------------------------------------------------------------

The observation is the walker configuration expressed relative to the query,
%
\begin{equation}
  s \;=\;
  \left(
    \frac{x_1 - x^{*}}{\Delta x}, \dots, \frac{x_N - x^{*}}{\Delta x},\;
    \frac{t_1 - t^{*}}{\Delta t}, \dots, \frac{t_N - t^{*}}{\Delta t},\;
    \frac{t^{*}}{\Delta t},\; \frac{x^{*}}{\Delta x}
  \right) \in \mathbb{R}^{2N+2}.
  \label{eq:state}
\end{equation}
%
Every entry is a displacement in grid cells. This is not cosmetic: the
quantities that govern admissibility --- two-sided coverage, collision,
conditioning --- all live at cell resolution and do not rescale with the
horizon. An earlier absolute-coordinate state carried no information about the
discretisation scale, and a fractional parameterisation $t/t^{*}$ was
considered and rejected for the mirror-image reason, since it makes ``one cell
from the query'' mean $0.02$ at one horizon and $0.1$ at another.

Two omissions are deliberate. The solution values $u_i$ are not observed, so
the policy chooses geometry without seeing the field it reconstructs. Nor is
$\mathcal{V}$ observed, even though it is $\mathcal{V}$ --- not the walker
positions --- that determines whether a stencil is admissible. The problem is
therefore partially observed, and the walker configuration is a summary
statistic rather than a state. We keep it because a fixed-size observation is
what allows one policy to serve queries at different locations;
\cref{sec:limitations} returns to the cost.

% ------------------------------------------------------------
\subsection{Actions and policy parameterisation}
\label{sec:mdp_action}
% ------------------------------------------------------------

A walker steps to a neighbouring cell,
%
\begin{equation}
  \mathcal{A} = \{-1,0,+1\} \times \{+1,+2\}
  \quad\text{in units of } (\Delta x, \Delta t),
  \label{eq:action_set}
\end{equation}
%
so evaluations always advance in time. Displacements with $\delta t = 0$ were
removed after they were observed to consume roughly half of all actions:
walkers advance monotonically and $\mathcal{V}$ already contains their own
trail, so a move at fixed $t$ almost always lands on an occupied cell. The
cost is recorded --- no walker can reposition laterally without also advancing
--- so lateral coverage of a time level must come from having several walkers
at different $x$.

A joint action selects both a walker and its move, giving a discrete space of
size $N \cdot |\mathcal{A}|$, decoded as
$i = \lfloor a/|\mathcal{A}| \rfloor$, $j = a \bmod |\mathcal{A}|$. The
flattened form matters. Factoring the policy into independent walker and
displacement heads,
$\pi(i,j \mid s) = \pi^{\mathrm{w}}(i \mid s)\,\pi^{\mathrm{m}}(j \mid s)$, is
a rank-one outer product: it has $(N-1) + (|\mathcal{A}|-1)$ free parameters
against $N|\mathcal{A}|-1$ for the full joint, and what lives in the missing
dimensions is exactly the correlation between which walker moves and where it
goes. A policy of the form ``if walker 3, go left; if walker 7, go right'' is
not expressible: mass on walkers $\{3,7\}$ and moves $\{\mathrm{left},
\mathrm{right}\}$ necessarily also puts mass on $(3,\mathrm{right})$ and
$(7,\mathrm{left})$.

We observed this as a gap between sampled and greedy performance. Under
deterministic rollout the action is assembled from two near-uniform marginals
and lands on a pair the policy never intended; under sampling the wrong pairs
are filtered by rejection and progress continues. A single softmax over the
joint set is full rank and closes the gap by construction.

Two earlier parameterisations failed for structural reasons. A fully shared
policy emitting a displacement for every walker at once has
$|\mathcal{A}|^{N}$ outputs and did not train: over $1.6\times10^{5}$ steps the
reached rate fell from $0.70$ to $0.55$. A policy applied to one walker at a
time reached the query reliably --- $0.04 \to 0.98$ over $10^{6}$ steps, with
the error improving from $1.1\times10^{-3}$ to $5.6\times10^{-4}$ --- but has
no mechanism for choosing \emph{which} walker to move, which is the
coordination the $N$ walkers exist to provide.

% ------------------------------------------------------------
\subsection{Admissibility masking}
\label{sec:mdp_mask}
% ------------------------------------------------------------

A proposal is masked if it leaves the domain or the ceiling
$t^{*} + R\,\Delta t$; if its target cell is already in $\mathcal{V}$; or if a
previous attempt there failed to admit a stencil, in which case the cell is
blocked until $\mathcal{V}$ has grown by a further $k$ points, on the reasoning
that the available stencil will have changed by then. The first two are
computable before acting; the third is not, since whether a solve succeeds is
only discovered by attempting it. Masking is therefore not a static function of
position --- it accumulates the episode's own history of failed solves.

Masking was the single most consequential implementation change, for a reason
that is a property of the MDP rather than of the learning algorithm. A rejected
proposal leaves \eqref{eq:state} unchanged, since the walker does not move. A
policy is a function of the observation, so having selected an action that is
rejected it is presented with the identical observation and selects the
identical action indefinitely; two-action cycles were also observed. Stochastic
sampling escapes such cycles only by wasting a large fraction of the budget.
Maskable PPO~\cite{huang2022masking} renormalises over admissible actions and
makes the pathology unreachable rather than merely improbable.

% ------------------------------------------------------------
\subsection{Termination}
\label{sec:mdp_term}
% ------------------------------------------------------------

A box of half-width $R$ cells is placed around the query,
$\mathcal{B}_R(z^{*})$, and the episode terminates when $z^{*}$ admits an
admissible stencil drawn from $\mathcal{B}_R(z^{*})$, or when a budget of
$K_{\max}$ steps is exhausted. No walker need land on $z^{*}$. Walkers may rise
above it, up to $t^{*} + R\,\Delta t$, so the box can be populated from both
sides in time; the final solve is correspondingly non-causal, whereas every
intermediate solve during the rollout uses causal neighbours only.

The stopping rule asks the same question as the scoring rule --- \emph{does an
admissible stencil exist at $z^{*}$} --- rather than the cheaper proxy
\emph{are $n$ evaluated points nearby}. The distinction is not cosmetic. Under
the proxy an episode can satisfy termination and then fail its final solve,
receiving the worst accuracy penalty regardless of anything the agent did; the
reward is then constant across those episodes and carries no gradient.

% ------------------------------------------------------------
\subsection{Reward}
\label{sec:mdp_reward}
% ------------------------------------------------------------

The return is terminal, and both branches are scored by the same expression. On
termination the operator is applied at $z^{*}$ and the error
$e = |u^{*} - \hat{u}^{*}|$ is compressed onto a bounded scale,
%
\begin{equation}
  e_{\mathrm{norm}}
  = \mathrm{clip}\!\left(\frac{\log_{10}(e/e_{\mathrm{tol}})}{3},\,0,\,1\right),
  \qquad
  r = -\,w_{\mathrm{err}}\,e_{\mathrm{norm}}
      \;-\; w_{\mathrm{time}}\,\frac{K}{K_{\max}},
  \label{eq:reward}
\end{equation}
%
with $e_{\mathrm{norm}} = 1$ when the final solve admits no stencil. We use
$\gamma = 1$: a discount would assert that evaluations placed early are worth
more than those placed late, which is not a preference we hold, since an
evaluation costs the same wherever it occurs.

\eqref{eq:reward} is the third reward we tried, and the first two failed in
ways that are worth reporting.

\paragraph{An unnormalised reward optimises the wrong term.}
With $r = -e - \lambda K$ and $\lambda = 10^{-3}$, the errors were of order
$10^{-5}$ while $\lambda K_{\max}$ was of order $1$, so the accuracy term
contributed nothing to the gradient. Over $3.1\times10^{6}$ steps training
drove the reached rate from $0.61$ to $1.00$ and the reward from $-0.615$ to
$-0.009$ --- essentially its maximum --- while the error did not move at all,
beginning at $2.04\times10^{-5}$ and ending at $2.06\times10^{-5}$. Episode
length collapsed from $225$ steps to $8.6$. The agent had optimised the reward
completely, by rushing.

\paragraph{Rescaling accuracy inverts the ordering.}
Introducing a weight $\beta$ on the accuracy term fixes that and breaks the
opposite thing. At $\lambda = 0.005$, $\beta = 100$ the reached rate fell from
$0.70$ to $0$ over $5\times10^{6}$ steps, every episode ran to the step budget,
and the explained variance collapsed from $0.70$ to $2\times10^{-4}$ --- while
the error on the episodes that still reached kept improving. Reaching had
become worse than timing out. We swept $\beta$ from $10^{2}$ to $10^{7}$ and
found no value satisfying $r_{\mathrm{timeout}} < r_{\mathrm{reach}}$ at every
horizon, because accumulated error eventually makes the reached branch worse
than the timeout branch and the only repair available by tuning is to shrink
the accuracy weight, that is, to declare poor accuracy acceptable.
Normalisation was adopted because the ordering has to hold structurally rather
than by calibration.

\paragraph{Shaping.}
A potential-based term densifies an otherwise terminal signal. With
$\phi(s) = -c_{\mathrm{shape}}\,N^{-1}\sum_i d_i/d_0$, where $d_i$ is the
normalised walker--query distance and $\phi \equiv 0$ at terminal states, the
contributions telescope to $-\phi(s^{(0)})$ and the optimal policy is
unchanged~\cite{ng1999policy}. An earlier version used $\min_i d_i$; under it a
move by any walker other than the current closest leaves $\phi$ unchanged, so
$N-1$ of the $N$ walker selections receive no shaping at all, and it rewards a
single walker sprinting ahead --- by \cref{sec:recon_bound} the geometry that
inflates $\|\mathbf{w}\|_1$.

% ------------------------------------------------------------
\subsection{Training}
\label{sec:mdp_training}
% ------------------------------------------------------------

We train with maskable PPO~\cite{schulman2017ppo,raffin2021sb3} on a
multilayer-perceptron policy over $12$ subprocess environments;
\cref{tab:config} lists the configuration. GPU execution was tested and found
slower than CPU: the policy network is small and the cost is dominated by the
environment step --- a neighbour search and a $3\times 5$ least-squares solve
--- so transfer overhead outweighs any gain.

The reconstruction environment enforces four admissibility conditions:
two-sided spatial coverage, a bound on $\operatorname{cond}(\mathbf{A})$, a
bound on the largest weight magnitude, and a lower bound on negative weights.
Carried unchanged into the RL setting all four together left the walkers barely
able to move: nearly every proposal was rejected, $K_{\max}$ had to be raised
substantially, and the reach rate collapsed. Under a sparse terminal reward
this is fatal in a specific way --- an agent that rarely reaches the query
never discovers that a reward exists there. We therefore disabled the two
weight constraints and retained the conditioning and coverage tests.

The environment accepts a schedule of queries drawn per episode, so that each
query supplies a distinct set of training examples relative to the state
\eqref{eq:state}. The implemented scheduler samples jointly and uniformly
rather than running a sequential curriculum.
\TODO{State plainly whether the featured run uses a single fixed query. If it
does, $t^{*}/\Delta t$ and $x^{*}/\Delta x$ are constant inputs and the
generalisation machinery is inactive --- say so here rather than letting a
reviewer find it.}
